\subsection{Two implementations, one mechanism}
\label{sec:impl}

Every sweep in this paper was run in PyTorch. The architecture is also implemented on an
autograd engine written from scratch for this project, and the two are compared here because
a mechanism that appears in only one implementation is a property of that implementation.

\paragraph{The engine is correct.} Against PyTorch at identical weights it agrees to
$4.4 \times 10^{-16}$ on the forward pass and $2.68 \times 10^{-15}$ on gradients, verified by
a finite-difference suite over every operator. It groks at four moduli ($113$, $119$, $121$
and $125$) at three seeds, and so reproduces the qualitative result of
Sections~\ref{sec:clock} and \ref{sec:strata} independently of PyTorch's autograd.

\paragraph{The two disagree about the magnitude of sparsity, and the cause is training
precision acting on PyTorch.} At byte-identical hyperparameters the engine's circuits read
$\mathrm{Gini}_{\text{mult}}$ about $0.2$ higher than PyTorch's at every modulus, $7$ of $7$
matched grokked pairs, with a smaller $\lVert W_{E} \rVert$ as well. The $2 \times 2$ resolves
it (Figure~\ref{fig:precision}):

\begin{center}
\begin{tabular}{lcc}
  \toprule
  & float32 & float64 \\
  \midrule
  from-scratch engine & $0.7671$ & $0.7536$ \\
  PyTorch & $\mathbf{0.5791}$ & $\mathbf{0.7728}$ \\
  \bottomrule
\end{tabular}
\end{center}

Each cell is tail $\mathrm{Gini}_{\text{mult}}$ at $n = 113$: the mean over the final
$10{,}000$ steps, never a single checkpoint, over three seeds of $40{,}000$ steps. The experiment
was pre-registered before either dtype was read. Every cell is recomputed from its artifact
rather than quoted.

Precision is inert in the engine and decisive in PyTorch. Take $\sigma_{\text{baseline}}$ to be
PyTorch at float32 and $\sigma_{\text{target}}$ the engine's float64 value. The effect size
\eqref{eq:effectsize} is then $f = +1.110$ on Gini and $f = +0.968$ on
$\lVert W_{E} \rVert$, both above the pre-registered threshold of $0.70$. The Gini value
exceeds $1$, meaning float64 moves PyTorch slightly \emph{past} the engine's position, and it
is given unclipped. Float64 also moves PyTorch's neuron tuning from $0.938$ to $0.770$. That
is both axes at once, onto the engine's joint position, which is what a single convergence
axis looks like. It dissolves an earlier worry that the two might be running different
algorithms.

\paragraph{This is an interaction, not a main effect, and the distinction is a retraction.}
We previously claimed that training precision changes measured sparsity, full stop, and tested
it by varying dtype \emph{inside the engine alone}. The effect size \eqref{eq:effectsize} came
back $f = -0.07$, the sign wrong and not merely the magnitude small, and the claim was
retracted. It stays retracted. What is true is the narrower sentence: training precision
changes the measured sparsity \emph{in PyTorch}.

Consequently every absolute sparsity number this paper reports from PyTorch is a float32 measurement,
and we name the regime beside each one. Comparisons between moduli, and comparisons within a
single arm, are unaffected; almost everything this paper claims is one of those. The
replication of \citet{nguyen2026dlogclock} in Section~\ref{sec:calib} matched their published
value because both arms are float32, and that agreement is real.

Before precision was reached, we eliminated the following candidates:
\begin{itemize}
  \item the measurement protocol;
  \item the hyperparameters, read from the kernel source rather than assumed;
  \item the tail window;
  \item the optimiser (the engine's AdamW is algebraically PyTorch's);
  \item gradients ($2.68 \times 10^{-15}$);
  \item minibatch noise (both are full-batch);
  \item the initialisation distribution;
  \item the train/test split (byte-identical, the same $3{,}830$ pairs);
  \item the execution substrate (PyTorch on CPU reads $0.5791$ against $0.5665$ on a T4, so
    no Kaggle number carries a hardware artefact);
  \item the specific initialisation draw.
\end{itemize}

\paragraph{What remains open is the engine's dtype-robustness, and it has no surviving
candidate.} The gap is explained; why the engine reaches the float64 answer at either dtype is
not. Under a pre-registration committed before any dtype was read we tested and eliminated
three mechanisms. They are residual type promotion in intermediate computations ($0$ of $115$
forward intermediates promoted, with a positive control at $100\%$), matmul accumulation order,
and softmax formulation (ratios $1.000$--$1.172$, inside the pre-registered equivalence band).
The question stays open with no candidate rather than with a proposed one.

% STE descriptive
\begin{figure}[t]
  \centering
  \includegraphics[width=0.62\textwidth]{../figures/F7_precision.pdf}
  \caption{\textbf{Training precision is an interaction, not a main effect.} Each bar is tail
    $\mathrm{Gini}_{\text{mult}}$ at $n = 113$ over three seeds and $40{,}000$ steps. The
    from-scratch engine reads essentially the same value at either dtype, and PyTorch moves
    by $0.19$. The experiment was pre-registered before either dtype was read. The
    figure regenerates from \texttt{src/viz/plots.py::precision\_2x2}.}
  \label{fig:precision}
\end{figure}
% STE end
